\documentclass[a4paper,DIV=16]{scrartcl}

\usepackage[english]{babel}
\usepackage{graphicx}
\usepackage[colorlinks=false]{hyperref}
\usepackage{tabularx}
\usepackage{placeins}

\author{Hugo Dischert}
\date{\today}
\title{PSL Week : Frugal AI}
\subtitle{Optimization of matrix multiplication}


\begin{document}
\maketitle
\tableofcontents
\FloatBarrier
\newpage
\section{General informations}
\paragraph{CPU specifications.}
The first part of this report deals with optimizing time and power consumption of matrix multiplication on CPU.
The CPU used is an Intel i7-12700, with 12 cores (8 P-Cores and 4 E-Cores) and 20 threads (P-Cores have hyperthreading).
It supports AVX2. 
Each P-Core has 48 KB of L1d cache and 1280 KB of L2 cache.
Each E-Core has 32 KB of L1d cache and all 4 share 2048 KB of L2 cache.
The L3 cache is 25 MB.

\paragraph{Methodology.}
Every numerical result is an average on at least three runs of the algorithm. 
Results found out to be always consistent across the three runs, with the last run being slightly longer than the first one during large matrix multiplication.
A small CPU overheat might be involved.

\paragraph{Power draw.}
The system power draw is 6W while idle and goes up to 25W during light use.
The monitor might be responsible for about 10 of those.
During heavy multithread computations, the system consumption rises to 170W.

\FloatBarrier
\section{Python implementation}
\paragraph{Naive implementation.}
A 3-loop Python implementation multiplies two 512x512 matrices in 29.95 seconds.
When replacing the inner loop with a numpy sum and vectorized multiplication, it achieves way faster results (0.63s for 512x512 matrices, 48 times faster).
This is because numpy operations are compiled.

\paragraph{JIT implementation.}
Using JIT to compile the two previous implementation, the multiplication become even faster (0.09s for the 3-loop implementation, 0.13s for the numpy one).
The numpy one is slower than the 3-loop one, probably because the compilation optimizes the inner loop as well as numpy sum and multiplication and it avoids calls to external functions.

\begin{figure}[h]
    \caption{Operations per second of each Python implementation, depending on the matrix size}
    \label{fig:Python FLOPS}
    \includegraphics[width=0.8\textwidth]{../saveq8/FLOPS.png}
\end{figure}
The curves of figure \ref{fig:Python FLOPS} tend to get closer as the matrix sizes grows.
The 3 top curves have dents on specific values.
Those values correspond to matrix sizes where it subsequently uses values bounded to the same cache slot, leading to many cache misses.
This disappear when we avoid power of 2 for the matrix size. (cf. figure \ref{fig:Without dents})

\begin{figure}[h]
    \caption{Python JIT implementation, where matrix sizes are multiple of 96. The dents disappears.}
    \label{fig:Without dents}
    \includegraphics[width=0.8\textwidth]{../saveq8/withoutdents.png}
\end{figure}

\paragraph{Numpy matrix multiplication.}
The numpy matrix multiplication is way faster than all 4 of these implementations as shown in figure \ref{fig:python FLOPSall}.
The speedup is about 100 for a matrix size of 1000x1000.

\begin{figure}[h]
    \caption{Previous 4 curves comparing to numpy implementation. The 3 top curves are multiple aliases for the same numpy feature.}
    \label{fig:python FLOPSall}
    \includegraphics[width=\textwidth]{../saveq8/FLOPSall.png}
\end{figure}

\FloatBarrier
\section{C implementation}
\paragraph{Naive 3-loop implementation.}
For a matrix size of 1920, the naive 3-loop implementation takes 17.88s without any flags.
The O2 flag achieves 5.02s and the O3 one 5.00s, which is a fiftieth of openBLAS 0.099s.
The figure \ref{fig:C naive} shows a huge performance drop when the matrix size reaches 1500, meaning the memory bandwidth is the performance bottleneck.

\begin{figure}[h]
    \caption{Performance of the 3-loop implementation.}
    \label{fig:C naive}
    \includegraphics[width=0.8\textwidth]{../saveC/naive03.png}
\end{figure}

\paragraph{Inner loop accumulator.}
By adding an accumulator, the memory is a little bit less accessed.
With the O2 flag, the multiplication of 1920x1920 matrices takes 4.53s.
The memory bandwidth is still the bottleneck, as shown in figure \ref{fig:accumulator}.

\begin{figure}[h]
    \caption{Performance of the 3-loop implementation with an accumulator.}
    \label{fig:accumulator}
    \includegraphics[width=0.8\textwidth]{../saveC/accumulator02.png}
\end{figure}

\paragraph{Transposed matrix.}
By transposing the first matrix, less memory bandwith is needed.
For a matrix of 1920x1920, it takes 2.87s with the O3 flag, achieving a thirtieth of openBLAS.
Adding the flags ffast-math, march=native and funroll-loops reduces the time to 0.48s, which is a fifth of openBLAS (tabular \ref{tab:transposed flags}).
The ffast-math flag seems to be the most important of those flags.
However, performances still degrade with the matrix size (cf. figure \ref{fig:transposed}).

\begin{figure}[h]
    \caption{Performance of the 3-loop implementation with transposition with different flags.}
    \label{tab:transposed flags}
    \begin{tabular}{c | c | c}
	Time in seconds & ffast-math & no ffast-math \\
	no march=native / no funroll-loops & 0.79 & 2.87 \\
	march=native / no funroll-loops & 0.56 & 2.90 \\
	no march=native / funroll-loops & 0.75 & 2.86 \\
	march=native / funroll-loops & 0.48 &  2.86
    \end{tabular}
\end{figure}


\begin{figure}[h]
    \caption{Performance of the 3-loop implementation with transposition and every flag.}
    \label{fig:transposed}
    \includegraphics[width=0.8\textwidth]{../saveC/transposed.png}
\end{figure}
\paragraph{Vectorization.}
With vectorization, 8 multiplications are done simultaneously.
This makes the code twice faster when using only O1, O2 or O3 flag.
However, when using ffast-math, the results are the same that without vectorization.
The compiler might have already done such optimization.
Furthermore, the memory bandwidth is still the bottleneck (cf. figure \ref{fig:Cvectorized}).

\begin{figure}[h]
    \caption{Performance of the 3-loop implementation with transposition and vectoization.}
    \label{fig:Cvectorized}
    \includegraphics[width=0.8\textwidth]{../saveC/Cvectorized.png}
\end{figure}

\paragraph{Kernel.}
Using a kernel avoids to reload the same data multiple times.
Best results were found with a kernel of 16x6 and using the flags O2, ffast-math, funroll-loops, march=native.
This kernel optimality comes from its small size: it entirely fits in the AVX2 registers.
For small matrix size, it achieves 70\% of openBLAS, but the figure \ref{fig:v5} illustrates a fast performance drop when the matrix size increase.


\begin{figure}[h]
    \caption{Performance of the kernel implementation.}
    \label{fig:v5}
    \includegraphics[width=0.8\textwidth]{../saveC/v5.png}
\end{figure}

\paragraph{Blocked matrix with kernel.}
The fastest blocks were those with $b_M = 64$, $b_N = 240$ and $b_K = 480$.
For a matrix size of 3840, it takes 1.07s while openBLAS takes 1.07s.
This implementation achieves 72\% of openBLAS.
Moreover, there is no performance drop when the matrix size increases anymore as shown in \ref{fig:v6}.
I don't have the explanation for the drop at 7680.

\begin{figure}[h]
    \caption{Performance of the blocked matrix with kernel implementation.}
    \label{fig:v6}
    \includegraphics[width=0.8\textwidth]{../saveC/v6_moyenne.png}
\end{figure}
\FloatBarrier
\section{Parallelization}
Increasing the number of threads increases the overall performance.
However, the performance is not linear in the number of threads (cf. figure \ref{fig:threads flops} and \ref{fig:threads speedup}).
OpenBLAS implementation has a huge drop going from 8 to 9 threads.
Since there are 8 P-cores on the test machine, openBLAS might poorly handle P and E cores simulteanously.


\begin{figure}[h]
    \caption{Flops of multithreaded blocked matrix with kernel implementation compared to openBLAS. Matrix size is 9600x9600.}
    \label{fig:threads flops}
    \includegraphics[width=0.8\textwidth]{../saveC/threads_flops9600.png}
\end{figure}
\begin{figure}[h]
    \caption{Speedup of multithreaded blocked matrix with kernel implementation. Matrix size is 9600x9600.}
    \label{fig:threads speedup}
    \includegraphics[width=0.8\textwidth]{../saveC/threads_speedup9600.png}
\end{figure}

\FloatBarrier
\section{GPU with CUDA}
In this section, the computation are done on an Nvidia T4 GPU.

\paragraph{cuBLAS implementation.}
The cuBLAS library multiplies matrices of size 4096 in 0.03 seconds on a T4 GPU, which is 2317 GFLOPS.
This is 30 times faster than openBLAS running on a single core on the test CPU and 15 times faster than openBLAS on 8 cores.


\begin{figure}[h]
    \caption{Number of operations per second on T4 GPU with cuBLAS depending on the matrix size}
    \label{fig:cuBLAS GFlops}
    \includegraphics[width=0.8\textwidth]{../saveGPU/cuBLAS_GFlops.png}
\end{figure}

On the figure \ref{fig:cuBLAS GFlops}, the GPU performance is beetween 1700 and 2300 GFlops, which is beetwen 12 and 16 times faster than multithreaded openBLAS.

The T4 GPU has 28 GFLOPS/W while the CPU (using openBLAS on 8 cores) has 6 GFLOPS/W.
The GPU is almost 5 times more efficient.

\begin{figure}[h]
    \caption{Number of operations per second on T4 GPU with naive CUDA implementation depending on the matrix size}
    \label{fig:cuda naive GFlops}
    \includegraphics[width=0.8\textwidth]{../saveGPU/cuda_naive_GFlops.png}
\end{figure}

\paragraph{Naive implementation.}
The CUDA naive implementation is really slow, with 19 GFLOPS for matrices of size 4096 (cf. figure \ref{fig:cuda naive GFlops}).
This is 3 times worse than single thread openBLAS and 20 times worse than openBLAS with 8 threads. 
It achieves only 1\% of cuBLAS.

\begin{figure}[h]
    \caption{Number of operations per second on T4 GPU with cuTLASS implementation depending on the matrix size}
    \label{fig:cutLASS}
    \includegraphics[width=0.8\textwidth]{../saveGPU/cuda_custom_GFlops.png}
\end{figure}

\paragraph{cuTLASS implementation.}
The custom cuTLASS implementation is around 1100 GFLOPS (cf. figure \ref{fig:cutLASS}), which is half the cuBLAS performance.

\begin{figure}[h]
    \caption{Number of operations per second on T4 GPU with FP16 precision depending on the matrix size}
    \label{fig:FP16}
    \includegraphics[width=0.8\textwidth]{../saveGPU/cuBLAS_FP16_GFlops.png}
\end{figure}

\paragraph{FP16 precision.}
When the precision is reduces to FP16, the matrix multiplication on T4 GPU achieves 17000 GFLOPS.
It is 8 times faster than FP32, which is consistent with the theoretical capacities of the T4 GPU being 8 times faster for FP16 than FP32.
It is still 4 times slower than the thoeoritical capabilities of this GPU.

\FloatBarrier
\section{CNN optimization}
In this section, all computations are done on a T4 GPU on Google Collab.

\paragraph{First augmentation of the dataset.}
I started by experimenting on different augmentations of the dataset, keeping the original backbone.
The options I used are the horizontal flip, the color jitter, the coarse dropout and the channel dropout.
Accuracy stayed between 78\% and 82\%.
Training the model with 50\% of gray-scaled images also achieves those performance, but with longer training steps.

\paragraph{Experimenting with the architecture.}
I implemented other architectures to achieve better accuracy.
First, I changed the network to take 128x128 images, while keeping the same layers.
It did not change the accuracy, the network was just slower.
To make it faster, I added an initial pull layer to convert to size 64x64.
It was predictible that the accuracy would be lower than just taking 64x64 as entries for the network, but I did not expect the only 56\% accuracy after training.

Then I used deeper neural networks.
With 5 convulutional layers, all with 3x3 feature size, and 3 sense layers, I achieved 86.91\% accuracy.

I also tried implemented a fully convulutional network : 5 layers of 3x3 convolution and 2x2 pooling, with feature size doubling at each layer, and an ending dense layer.
I also tried a fully dense network, with 3 dense layers.
It was really slow to train and had poor accuracy (60 5\%).

\paragraph{Normalization layers.}
I went back to the deep architecture with 5 convolutional layers and 3 dense ones, and added normalization layers between each one. It only achieves 84\% accuracy.

A fully convulutional network with a single convulutional layer before the first pool and then 3 convulutional layers and a normalization layer between each pool (4 times) achieved 90.9\% accuracy but found out to be a little bit too slow.
Replacing the first set of 3 convolutional layers with a single one was still too slow.
To gain a little bit efficiency, I inverted the first normalizing and pulling layers.
The network was now running in about 295 ms, and achieved 89.26\%.
Adding data augmentation, the accuracy rised to 90.58\%.

I tried augmenting the size of the dense layers, but i went too far : my network had 9 million parameters.
The network overfitted the data and classified all test images as dogs.

I ended by returning to 64x64 images. I tried the same things which all achieved around 90\%.


\end{document}
